% Generated draft: requires mathematical and source review.
\documentclass[11pt]{article}
\usepackage[margin=25mm]{geometry}
\usepackage{fontspec}
\setmainfont{Latin Modern Roman}
\usepackage{amsmath,amssymb,mathtools,unicode-math}
\setmathfont{Latin Modern Math}
\usepackage{enumitem,needspace,xurl,hyperref}
\setlength{\parindent}{0pt}
\setlength{\parskip}{5pt}
\setlength{\emergencystretch}{3em}
\newcommand{\N}{\mathbb{N}}
\newcommand{\Z}{\mathbb{Z}}
\newcommand{\Q}{\mathbb{Q}}
\newcommand{\R}{\mathbb{R}}
\newcommand{\C}{\mathbb{C}}
\newcommand{\F}{\mathbb{F}}
\newcommand{\A}{\mathbb{A}}
\newcommand{\PP}{\mathbb{P}}
\newcommand{\E}{\mathbb{E}}
\newcommand{\eps}{\varepsilon}
\newcommand{\dd}{\,\mathrm{d}}
\begin{document}
\section*{UnsolvedMath ID: 30002379}
\textbf{Source problem code:} catalogue entry
\subsection{Definitions and mathematical statement}
Let \(p>37\) be a prime. The supplied record does not define its modular-curve notation. Under the usual conventions, \(X_0(p)\) is the smooth projective modular curve over \(\Q\) whose noncuspidal geometric points classify pairs \((E,C)\), where \(E\) is an elliptic curve and \(C\subset E[p]\) is a cyclic subgroup of order \(p\). The Atkin--Lehner involution is
\[
w_p(E,C)=(E/C,E[p]/C),
\]
and \(X_0^+(p)=X_0(p)/\langle w_p\rangle\).

For the other two curves, let
\[
C_{\mathrm{sp}}=\{\operatorname{diag}(a,b):a,b\in\F_p^\times\}
\subseteq\operatorname{GL}_2(\F_p).
\]
Choose an \(\F_p\)-basis of \(\F_{p^2}\), and let \(C_{\mathrm{ns}}\subseteq\operatorname{GL}_2(\F_p)\) be the image of \(\F_{p^2}^{\times}\) acting on \(\F_{p^2}\) by multiplication. For \(\star\in\{\mathrm{sp},\mathrm{ns}\}\), put
\[
N_\star=\{g\in\operatorname{GL}_2(\F_p):gC_\star g^{-1}=C_\star\}.
\]
Under the standard normalizer convention, \(X_\star^+(p)\) is the smooth projective modular curve over \(\Q\) associated with \(N_\star\). Its noncuspidal geometric points classify elliptic curves with a basis of their \(p\)-torsion modulo the action of \(N_\star\), up to isomorphism. Changing the chosen basis used to define \(C_{\mathrm{ns}}\) gives an isomorphic modular curve.

For any of these curves \(X\), write \(\operatorname{Cusps}(X)\) for its geometric cusps. A noncuspidal geometric point is a \emph{CM point} if it has a modular representative with underlying elliptic curve \(E/\overline{\Q}\) satisfying
\[
\operatorname{End}_{\overline{\Q}}(E)\supsetneq\Z.
\]
On \(X_0^+(p)\), this means the image of a CM point of \(X_0(p)\); the condition is independent of the representative because the two elliptic curves are isogenous. Denote the resulting set by \(\operatorname{CM}(X)\). Rationality here means rationality of the point on the coarse modular curve; it does not require a chosen elliptic-curve representative or level structure to be defined over \(\Q\).

With these explicitly stated conventions, the supplied question is whether
\[
\forall p\text{ prime},\quad p>37,\quad
\forall X\in\{X_0^+(p),X_{\mathrm{sp}}^+(p),X_{\mathrm{ns}}^+(p)\},\qquad
X(\Q)\subseteq\operatorname{Cusps}(X)\cup\operatorname{CM}(X).
\]
Thus every rational noncuspidal point is required to be a CM point. This is a proposed assertion, not an established conclusion here.

There is an unresolved qualification: the supplied literature assessment calls the first case the Borel case, whereas the usual Borel modular curve is \(X_0(p)\), and the displayed source statement uses its Atkin--Lehner quotient \(X_0^+(p)\). These are different rational-point questions. The draft retains the supplied superscript \(+\); the intended first curve and the applicability of the supplied literature assessment require source review.
\subsection{Short English statement}
For every prime \(p>37\), the supplied question asks whether every rational point on each of the three displayed modular curves is either a cusp or a CM point. The notation \(X_0^+(p)\) is retained, but its compatibility with the source's description of a Borel case remains unresolved. The supplied partially solved label is dated curation, not a conclusion independently established here.
\subsection{Sources}
\begin{itemize}
\item Supplied local UnsolvedMath registry record, permanent ID \(30002379\), exact source problem code \texttt{catalogue entry}, title \emph{Rational Points on Prime-Level Modular Curves}, in the Oberwolfach Reports Open Problems collection. The proposed year is recorded as 2013; no individual proposer is supplied.
\item Supplied source citation: \emph{Explicit Methods in Number Theory (2014)}, Oberwolfach Reports, DOI \texttt{10.4171/owr/2013/35}, pp.\ 2035--2082. \url{https://doi.org/10.4171/owr/2013/35}. The literature list labels this same reference \emph{Oberwolfach Report 35/2013}; the supplied year descriptions are retained without reconciliation.
\item Supplied literature reference: \emph{Serre's uniformity problem in the split Cartan case}. \url{https://annals.math.princeton.edu/2011/173-1/p13}.
\item Supplied literature reference: \emph{Serre's uniformity question and proper subgroups of the nonsplit Cartan}. \url{https://arxiv.org/abs/2305.17780}.
\item The supplied literature assessment was checked on \texttt{2026-08-21}. It labels the combined problem partially solved, reports the Borel and split-Cartan rational-point problems as solved in the required large-prime range, and describes the corresponding nonsplit-Cartan problem as a remaining central obstruction to uniformity. This is supplied dated curation, not independently verified; its claims and their applicability to the displayed curves require review.
\item The registry records statement status \texttt{corrected\_verified} and says that the clean statement was checked against the official Oberwolfach report. This is a supplied verification claim, not verification performed for this draft.
\end{itemize}
\end{document}
